\documentclass[11pt]{article}
\usepackage[margin=1in]{geometry}
\usepackage{amsmath,amssymb,amsthm}
\usepackage{array,longtable,booktabs}
\usepackage[hidelinks]{hyperref}

\newtheorem{theorem}{Theorem}
\newtheorem{proposition}[theorem]{Proposition}
\newtheorem{corollary}[theorem]{Corollary}
\newtheorem{example}[theorem]{Example}
\newcommand{\Ann}{\operatorname{Ann}}
\newcommand{\Char}{\operatorname{Char}}

\title{P6-T03: One-Ray Characteristic Data Do Not Select a Lorentzian Conformal Class}
\author{The AEther-Flow Research Project}
\date{2026-07-26}

\begin{document}
\maketitle

\begin{center}
\fbox{\parbox{0.92\textwidth}{
\textbf{Status and authority.}
This is a task-local \emph{draft/control} result about the proposal-only
P6-T02 source-extension candidate.  It is a precise candidate-scoped
obstruction, not a canonical-ontology edit, source-law adoption, physical
causality result, effective-metric construction, Gate Chair verdict, global
no-go theorem, publication, or completed derivation.
}}
\end{center}

\section{Question and fixed source input}

Let \(S\) be the smooth four-dimensional source arena carried as explicit
primitive debt in the current project, and let \(V\) be the nowhere-zero,
positively oriented source vector field introduced only by the P6-T02
proposal.  The linearized principal polynomial of that proposal is
\[
   P_x(\xi)=\xi(V_x).
\]
Consequently,
\[
   \Char_x^*=\Ann(V_x)\setminus\{0\},\qquad
   \Gamma_x=\{\xi\in T_x^*S:\xi(V_x)>0\},\qquad
   \Gamma_x^\circ=\{rV_x:r\geq 0\}.
\]
The source-order datum is the forward \(V\)-flow reachability preorder.
Denote this collection by
\[
   \mathcal D_1(S,V)
   =\bigl([V]_+,\Ann(V),\preceq_V\bigr).
\]
Here \([V]_+\) is an oriented ray, not a Lorentzian null cone.  A putative
conformal reconstructor would have to define a source-natural assignment
\[
   \mathfrak C:\mathcal D_1(S,V)\longmapsto[g],
\]
where \([g]\) is a nondegenerate Lorentzian conformal class, without importing
a target atlas, target metric, general-relativistic null cone, physical clock,
detector response, or benchmark behavior.

\section{Assumption table}

\begin{longtable}{>{\raggedright\arraybackslash}p{0.12\textwidth}
                  >{\raggedright\arraybackslash}p{0.34\textwidth}
                  >{\raggedright\arraybackslash}p{0.43\textwidth}}
\toprule
ID & Assumption & Status and consequence\\
\midrule
\endhead
A1 & \(S\) is a smooth four-dimensional source manifold. &
Inherited as explicit primitive debt.  It is not identified with physical
spacetime.\\
A2 & \(V\) is smooth, nowhere zero, and has a chosen positive orientation. &
Proposal-only P6-T02 source-extension data.  Current ontology does not derive
or select it.\\
A3 & Only \([V]_+\), \(\Ann(V)\), and the \(V\)-flow preorder are available. &
This is the exact one-ray characteristic information used by the
reconstruction test.\\
A4 & For the strongest countermodel, \(V\) is additionally stipulated to be
future timelike for every comparison metric. &
This favorable causal stipulation still does not select a transverse
conformal shape.  It is a test assumption, not a physical conclusion.\\
A5 & No Lorentzian scale is fixed. &
P6-T03 addresses conformal selection only; clock, rod, volume, and scale
calibration remain open.\\
A6 & The historical scoped source-extension \(g_{\rm eff}\) object is not
used as a premise. &
Its protected scope is preserved and is not expanded into an unscoped
physical metric.\\
A7 & No target-side geometric data are imported. &
The result is source-side and candidate-scoped.\\
\bottomrule
\end{longtable}

\section{Exact mismatch with a Lorentzian null cone}

\begin{theorem}[Hyperplane/null-cone mismatch]
\label{thm:hyperplane}
At any \(x\in S\) with \(V_x\neq0\), the three-dimensional linear hyperplane
\(\Ann(V_x)\subset T_x^*S\) cannot equal the null cone of any nondegenerate
Lorentzian inverse bilinear form on \(T_x^*S\).
\end{theorem}

\begin{proof}
Assume that a Lorentzian inverse form \(q\) vanishes on every covector in
\(\Ann(V_x)\).  Because \(\Ann(V_x)\) is a linear subspace, for any
\(\alpha,\beta\in\Ann(V_x)\) the polarization identity gives
\[
  2q(\alpha,\beta)
  =q(\alpha+\beta,\alpha+\beta)-q(\alpha,\alpha)-q(\beta,\beta)=0.
\]
Thus \(\Ann(V_x)\) would be a three-dimensional totally isotropic subspace.
A nondegenerate real bilinear form of Lorentzian signature has Witt index one,
so every totally isotropic subspace has dimension at most one.  This is a
contradiction.  Equivalently, a Lorentzian null cone is quadratic and
nonlinear, whereas \(\Ann(V_x)\) is a linear hyperplane.
\end{proof}

\begin{corollary}
The equation \(P_x(\xi)=\xi(V_x)\) does not itself provide the quadratic
principal polynomial whose zero locus could determine a Lorentzian conformal
class.  Treating \(\Ann(V_x)\) as such a null cone would import structure not
contained in P6-T02.
\end{corollary}

\section{Infinite conformal nonselection family}

\begin{theorem}[One-ray conformal nonselection]
\label{thm:nonselection}
Even after stipulating that the oriented ray \([V]_+\) is future timelike, the
data \(\mathcal D_1(S,V)\) do not determine a Lorentzian conformal class.
Locally, infinitely many pairwise nonconformal Lorentzian metrics realize the
same \([V]_+\), \(\Ann(V)\), and forward \(V\)-flow preorder.
\end{theorem}

\begin{proof}
By the flow-box theorem, choose local source coordinates
\((s,y^1,y^2,y^3)\) with \(V=\partial_s\).  Let \(H=(H_{ij})\) be any smooth
positive-definite symmetric \(3\times3\) matrix field and define
\[
   g_H=-ds^2+H_{ij}\,dy^i\,dy^j.
\]
Each \(g_H\) is Lorentzian, \(g_H(V,V)=-1\), and it has the same oriented
timelike ray \([V]_+\).  The covector hyperplane \(\Ann(V)\) and the forward
flow preorder depend only on \(V\), so they also agree for every \(H\).

Take
\[
   H_0=\operatorname{diag}(1,1,1),\qquad
   H_1=\operatorname{diag}(2,1,1).
\]
If \(g_{H_1}=\Omega^2g_{H_0}\), evaluation on \((V,V)\) gives
\(\Omega^2=1\).  Evaluation on
\((\partial_{y^1},\partial_{y^1})\) would then require \(2=1\), a
contradiction.  Hence the two metrics are not conformally related.
Continuous positive-definite families of \(H\) give infinitely many such
classes.
\end{proof}

\begin{example}[Different null cones with identical P6-T02 data]
For \(H_0\), the vector
\(\partial_s+\partial_{y^1}\) is null.  For \(H_1\), its squared norm is
\(-1+2=1\).  Thus the null cones differ even though the P6-T02 ray,
annihilator, and preorder agree exactly.
\end{example}

\begin{example}[A chosen null ray is also insufficient]
If one instead stipulates that \(V=\partial_s\) is null, the family
\[
   \widetilde g_a=2\,ds\,du+a^2(dy^1)^2+(dy^2)^2,\qquad a>0,
\]
is Lorentzian and shares the null ray \([V]\).  Distinct constant values of
\(a\) are not conformally related once the \(ds\,du\) coefficient is fixed.
One selected null generator therefore does not determine a full null cone.
\end{example}

\section{Uniqueness, degeneracy, and topology controls}

\begin{enumerate}
\item \textbf{Nonzero one-ray control.}  The family \(g_H\) proves
nonuniqueness even under the favorable timelike stipulation and smooth
positive-definite transverse geometry.
\item \textbf{Zero-field degeneracy.}  If \(V_x=0\), then
\(P_x\equiv0\), \(\Ann(V_x)=T_x^*S\), and neither a hyperbolicity direction nor
a distinguished propagation ray survives.
\item \textbf{Disconnected source control.}  On distinct connected
components, independent choices of \(H\) preserve the same componentwise
one-ray data; no cross-component conformal compatibility rule follows.
\item \textbf{Multiple-sector control.}  Two noncollinear sector rays may
improve directional information, but a finite set of rays still does not by
itself supply a smooth nondegenerate quadratic cone, signature law,
universality rule, or uniqueness theorem.
\item \textbf{Scale control.}  Even a successfully selected conformal class
would leave multiplication by a positive scalar field open.  No clock, rod,
volume, or duration calibration is introduced here.
\end{enumerate}

\section{Precise obstruction and reopening criterion}

\begin{proposition}[Candidate-scoped obstruction]
There is no well-defined conformal reconstruction map
\(\mathfrak C\) from only the P6-T02 datum \(\mathcal D_1(S,V)\): identical
input data admit nonconformal outputs.  The obstruction scope is the
proposal-only one-ray candidate and any reconstruction rule whose complete
input is only \([V]_+\), \(\Ann(V)\), and \(\preceq_V\).
\end{proposition}

This is not a theorem that the project can never derive effective conformal
geometry.  Same-milestone continuation remains open if it supplies materially
new source-side input, for example:
\begin{itemize}
\item a source-derived smooth, nondegenerate Lorentzian cone or quadratic-form
law;
\item a source-derived transverse positive-definite conformal shape together
with compatibility, covariance, and universality selectors; or
\item a distinct theorem showing that richer multi-sector characteristic data
uniquely determine a conformal class under explicit source-derived
hypotheses.
\end{itemize}
Any such input remains \emph{draft/control}, \emph{proposal-only}, and
\emph{source-extension data} until a separate protected human decision.
Current adoption is therefore
\texttt{blocked\_adoption\_open\_continuation}.

\section{Distance-to-GR disposition}

P6-T03 is mathematically decisive within its stated scope: it proves exact
hyperplane mismatch and conformal nonselection for the P6-T02 input.  It
narrows the \(g_{\rm eff}\) burden by identifying the missing conformal-shape
primitive, but it does not construct or adopt \(g_{\rm eff}\), discharge
metric scale, establish physical causality, couple matter, derive Einstein
equations, or promote the exact-GR benchmark.

\section{Project source note}

The fixed source basis is the registered P6-T02 transport candidate, the
canonical target-side geometry dictionary, the historical scoped
source-extension \(g_{\rm eff}\) candidate and Gate Chair record, the v21
implementation plan, and the Distance-to-GR burden map.  Those project
sources are held at the exact hashes recorded in the accompanying
machine-readable specification.  The linear-algebra and counterexample
arguments above are proved directly and use no external theorem beyond
standard finite-dimensional bilinear-form facts reproduced in the proof.

\end{document}
